\chapter{Delta X2 G-code Protocol Reference}
\label{app:gcode}
\index{G-code}\index{protocol!serial}%

This appendix is the authoritative description of the G-code dialect and
communication protocol of the Delta X2 robot as used by this project.

\section{Communication protocol}
\begin{table}[h]
  \centering
  \begin{tabularx}{\textwidth}{lD}
    \toprule
    \textbf{Parameter} & \textbf{Value} \\
    \midrule
    Transport & USB serial (CDC), typically \cfg{/dev/ttyUSB*} or \cfg{/dev/ttyACM*} \\
    Baud rate & 115\,200 \\
    Line ending & \cfg{\textbackslash n} (newline) \\
    Handshaking & The robot answers commands; for critical synchronisation
      use the \cfg{FEEDBACK} parameter (section~\ref{sec:gcode-feedback}). \\
    \bottomrule
  \end{tabularx}
  \caption{Serial link parameters}
\end{table}

\section{Identification and status}

\subsection{\texttt{IsDelta}}
\index{G-code!IsDelta@\texttt{IsDelta}}\index{handshake}%
Queries the device to confirm it is a Delta X robot.
\begin{itemize}
  \item \textbf{Response:} \cfg{YesDelta}
\end{itemize}
The application performs this handshake on every connection and refuses
to proceed if the answer does not arrive.

\subsection{\texttt{DeltaState}}
\index{G-code!DeltaState@\texttt{DeltaState}}%
Queries the current operating state of the robot.

\begin{table}[h]
  \centering
  \begin{tabularx}{\textwidth}{lD}
    \toprule
    \textbf{Response} & \textbf{Meaning} \\
    \midrule
    \cfg{Free} & Idle and ready. \\
    \cfg{Running} & Executing a command. \\
    \cfg{Wait} & Waiting for input or synchronisation. \\
    \cfg{Almostdone} & Near the end of a movement. \\
    \cfg{Done} & Movement or command completed. \\
    \bottomrule
  \end{tabularx}
  \caption{\cfg{DeltaState} responses}
\end{table}

\subsection{\texttt{Position}}
\index{G-code!Position@\texttt{Position}}\index{position feedback}%
Returns the current coordinates of the robot.
\begin{itemize}
  \item \textbf{Response format:} \cfg{X:<val> Y:<val> Z:<val> W:<val> U:<val>}
\end{itemize}

\begin{notebox}
  \cfg{Position} is not used by the application yet; it is the basis for
  the planned belt-motion compensation (robot position feedback for the
  trajectory planner --- see \cfg{documentation/TODO.md}).
\end{notebox}

\section{Movement commands}

\subsection{\texttt{G00} / \texttt{G01} --- linear movement}
\index{G-code!G01@\texttt{G01}}\index{feed rate}\index{movement}%
Moves the end-effector to the specified coordinates.
\begin{itemize}
  \item \textbf{Syntax:} \cfg{G01 X<val> Y<val> Z<val> F<speed>}
  \item \cfg{F} is the feed rate in mm/min (configured via
    \cfg{robot.feed\_rate}).
\end{itemize}
\begin{lstlisting}
G01 X10.5 Y-20.0 Z-150.0 F2000
\end{lstlisting}

\subsection{\texttt{G28} --- homing}
\index{G-code!G28@\texttt{G28}}\index{homing}%
Moves all axes to their home position (top endstops).
\begin{itemize}
  \item \textbf{Origin (X0 Y0 Z0):} after homing, the robot is positioned
    at the top centre.
  \item \textbf{Z axis:} moves downward (negative values). Z$=0$ is the
    top; Z$=-200$ is near the bottom of the workspace.
\end{itemize}

\subsection{\texttt{G90} / \texttt{G91} --- positioning mode}
\index{G-code!G90@\texttt{G90}}\index{G-code!G91@\texttt{G91}}%
\begin{itemize}
  \item \cfg{G90}: interpret subsequent coordinates as absolute positions
    relative to the origin. The application sets this mode right after
    homing.
  \item \cfg{G91}: interpret subsequent coordinates as displacements from
    the current position (useful for jogging).
\end{itemize}

\section{Synchronisation --- \texttt{FEEDBACK}}
\label{sec:gcode-feedback}
\index{G-code!FEEDBACK@\texttt{FEEDBACK}}\index{synchronisation}%
\cfg{FEEDBACK:<string>} can be appended to any G-code command. The robot
echoes \cfg{<string>} back on the serial port once the command has been
\emph{physically executed} --- not merely received.

\begin{lstlisting}
G01 X0 Y0 Z-50 FEEDBACK:done
\end{lstlisting}
Response (after the move finishes):
\begin{lstlisting}
done
\end{lstlisting}

\begin{notebox}
  The robot driver appends a unique \cfg{FEEDBACK:sync\_<n>} id to every
  command and blocks until the echo returns. This is what makes a
  completed \cfg{move\_to} call mean ``the arm has arrived'', and it is
  why every command has an overall timeout: a lost echo surfaces as an
  error instead of a hang.
\end{notebox}

\section{End effector controls}
\index{G-code!M03@\texttt{M03}}\index{G-code!M05@\texttt{M05}}\index{G-code!M112@\texttt{M112}}\index{gripper}\index{emergency stop}%
\begin{table}[h]
  \centering
  \begin{tabularx}{\textwidth}{lD}
    \toprule
    \textbf{Command} & \textbf{Effect} \\
    \midrule
    \cfg{M03} & Output on (open gripper / start vacuum). \\
    \cfg{M05} & Output off. \\
    \cfg{M112} & Emergency stop (used by the E-stop path). \\
    \bottomrule
  \end{tabularx}
  \caption{M-codes}
\end{table}

\section{Physical workspace (SP-X2)}
\index{workspace!physical limits}\index{SP-X2}%
\begin{table}[h]
  \centering
  \begin{tabularx}{\textwidth}{lD}
    \toprule
    \textbf{Axis} & \textbf{Range} \\
    \midrule
    X & $-160$\,mm to $+160$\,mm \\
    Y & $-160$\,mm to $+160$\,mm \\
    Z & $-200$\,mm to $0$\,mm ($0$ at the top) \\
    Working diameter & 320\,mm \\
    \bottomrule
  \end{tabularx}
  \caption{Workspace of the SP-X2 model}
\end{table}

These are the hard physical limits; the configured workspace
(\cfg{[robot]} section, chapter on configuration) must always be equal to
or smaller than them.

\section{Quick reference}
\begin{table}[h]
  \centering
  \begin{tabularx}{\textwidth}{lD}
    \toprule
    \textbf{Command} & \textbf{Meaning} \\
    \midrule
    \cfg{IsDelta} & Handshake; the robot answers \cfg{YesDelta}. \\
    \cfg{DeltaState} & Operating state query. \\
    \cfg{Position} & Current coordinates query. \\
    \cfg{G28} & Home all axes (to the top; X0 Y0 Z0 afterwards). \\
    \cfg{G90} / \cfg{G91} & Absolute / relative positioning mode. \\
    \cfg{G01 X.. Y.. Z.. F..} & Linear move at feed rate \cfg{F} (mm/min). \\
    \cfg{M03} / \cfg{M05} & End effector (gripper/pump) on / off. \\
    \cfg{M112} & Emergency stop. \\
    \cfg{FEEDBACK:<id>} & Echoed back when the command has physically
      completed. \\
    \bottomrule
  \end{tabularx}
  \caption{Protocol summary}
\end{table}
